%==============================================================================
% Chapitre 2 - Cinématique
%==============================================================================

\section{Introduction}

La cinématique est la branche de la mécanique qui décrit les mouvements sans
chercher à expliquer leurs causes.

Elle s'intéresse notamment à :

\begin{itemize}
    \item la position ;
    \item la vitesse ;
    \item l'accélération ;
    \item la trajectoire.
\end{itemize}

Contrairement à la dynamique, la cinématique ne s'intéresse pas aux forces.

Elle cherche uniquement à répondre à la question :

\begin{quote}
Comment un objet se déplace-t-il ?
\end{quote}

\section{Pourquoi la cinématique est importante en balistique}

La balistique extérieure consiste essentiellement à étudier le mouvement d'un
projectile entre son départ du canon et son impact.

Pour comprendre ce mouvement, il est nécessaire de savoir décrire :

\begin{itemize}
    \item où se trouve le projectile ;
    \item à quelle vitesse il se déplace ;
    \item comment sa vitesse évolue ;
    \item quelle trajectoire il suit.
\end{itemize}

Ces notions constituent le vocabulaire de base de toute étude balistique.

\section{Position}

La position permet de localiser un objet dans l'espace.

Pour cela, il est nécessaire de définir un repère.

Dans le cadre de ce manuel, nous utiliserons généralement un repère
cartésien constitué de trois axes :

\begin{itemize}
    \item l'axe X ;
    \item l'axe Y ;
    \item l'axe Z.
\end{itemize}

La position d'un objet peut alors être définie par trois coordonnées.

\section{Trajectoire}

La trajectoire est l'ensemble des positions successivement occupées par un
objet au cours du temps.

Selon les conditions de mouvement, une trajectoire peut être :

\begin{itemize}
    \item rectiligne ;
    \item circulaire ;
    \item parabolique ;
    \item quelconque.
\end{itemize}

En balistique réelle, la trajectoire d'un projectile n'est généralement pas
une parabole parfaite.

La gravité, la traînée aérodynamique et d'autres phénomènes modifient sa
forme.

\section{Distance parcourue et déplacement}

Deux notions sont souvent confondues :

\begin{itemize}
    \item la distance parcourue ;
    \item le déplacement.
\end{itemize}

La distance parcourue correspond à la longueur totale du trajet effectué.

Le déplacement correspond à la différence entre la position initiale et la
position finale.

Un objet peut parcourir plusieurs kilomètres et revenir à son point de départ.

Dans ce cas :

\begin{itemize}
    \item la distance parcourue n'est pas nulle ;
    \item le déplacement est nul.
\end{itemize}

Cette distinction deviendra importante lors de l'étude des trajectoires.

\section{Vitesse}

La vitesse décrit la variation de position au cours du temps.

Dans le langage courant, on utilise souvent uniquement sa valeur.

Par exemple :

\begin{itemize}
    \item 10 km/h ;
    \item 50 km/h ;
    \item 850 m/s.
\end{itemize}

Cependant, en physique, la vitesse possède également une direction.

Deux objets peuvent avoir la même vitesse en valeur tout en se déplaçant dans
des directions différentes.

\section{Unités utilisées}

Dans ce manuel, les vitesses sont généralement exprimées en mètres par
seconde (m/s).

Certaines vitesses peuvent également être données en kilomètres par heure
(km/h) lorsqu'elles correspondent à des usages courants.

La conversion est :

\begin{equation}
1~\mathrm{m/s}=3,6~\mathrm{km/h}
\end{equation}

\section{Exemple : projectile de MAG58}

Une balle quittant le canon d'une MAG58 avec une vitesse initiale de
850~m/s possède :

\begin{itemize}
    \item une vitesse très élevée ;
    \item une direction initiale déterminée par l'orientation du canon.
\end{itemize}

Pendant son vol :

\begin{itemize}
    \item sa vitesse diminue sous l'effet de la traînée ;
    \item sa direction évolue sous l'effet de la gravité.
\end{itemize}

\section{Accélération}

L'accélération mesure l'évolution de la vitesse au cours du temps.

Une accélération peut correspondre :

\begin{itemize}
    \item à une augmentation de la vitesse ;
    \item à une diminution de la vitesse ;
    \item à un changement de direction.
\end{itemize}

Ainsi, un objet se déplaçant à vitesse constante sur une trajectoire
circulaire possède une accélération.

Sa vitesse conserve la même valeur mais sa direction change continuellement.

\section{Mouvements fondamentaux}

Les principaux types de mouvements étudiés en mécanique sont :

\begin{itemize}
    \item le mouvement rectiligne uniforme ;
    \item le mouvement rectiligne uniformément accéléré ;
    \item le mouvement circulaire ;
    \item le mouvement curviligne.
\end{itemize}

Le mouvement d'un projectile réel appartient généralement à la dernière
catégorie.

\section{Application à la balistique}

La cinématique permet de décrire :

\begin{itemize}
    \item la trajectoire d'un projectile ;
    \item le temps de vol ;
    \item la vitesse résiduelle ;
    \item la chute ;
    \item les déplacements d'une cible ;
    \item les mouvements d'une plateforme de tir.
\end{itemize}

Elle constitue donc la base mathématique nécessaire à l'étude de la
balistique extérieure.

%\section{À retenir}

\begin{aretenir}

\begin{itemize}
    \item La cinématique décrit les mouvements sans étudier leurs causes.
    \item La position permet de localiser un objet dans l'espace.
    \item La trajectoire est l'ensemble des positions successives d'un objet.
    \item La vitesse possède une valeur et une direction.
    \item L'accélération traduit une variation de vitesse ou de direction.
    \item La cinématique constitue l'un des fondements de la balistique.
\end{itemize}

\section{À retenir}

\end{aretenir}

\begin{rigueur}
    \begin{equation}
    v = \frac{dx}{dt}
    \end{equation}

    \begin{equation}
    a = \frac{dv}{dt}
    \end{equation}
\end{rigueur}
